\section{Utility Evaluation of Downstream Tasks}\label{sec:utility}
Most market generator research focuses on statistical similarity while
overlooking practical downstream effects \cite{Horvath2025, Stenger24}.
While some work explores option hedging \cite{Tanaka25,Boursin22}, benchmarks
like TSGBench lack broader utility experiments. We introduce three
downstream tasks \ref{contribute:2} under a train-synthetic-test-real (TSTR)
framework \cite{Leznik21, Stenger24} --- option delta hedging (options
market maker), portfolio hedging (equity market maker), and alpha generation
(hedge fund manager) --- demonstrating synthetic data utility in realistic
settings \cite{Horvath2025}.

\subsection{Overview of Utility Evaluation}
    All three tasks follow a common protocol. Each task
    \(\mathcal{T}\) is specified by a loss function
    \(\mathcal{L}_{\mathcal{T}}\) over the model's policy \(\theta\), and
    the optimal policy is
    \begin{equation}
      \theta^{*} \coloneq \underset{\theta}{\arg\min}\;
      \mathcal{L}_{\mathcal{T}}(\theta),
      \label{eq:utility-optimal-policy}
    \end{equation}
    (or \(\arg\max\) for reward-type objectives). Policies are trained
    under the TSTR framework of §\ref{sec:utility-training-framework} on
    the data partition of §\ref{subsec:data-preprocess} with rolling-window
    validation for hyperparameter tuning \cite{Leznik21}, and scored with
    the metric toolbox of §\ref{sec:utility-evaluation}.

    \subsubsection{Utility Training and Evaluation Framework}
    \label{sec:utility-training-framework}
      Consider a market generator \(G_g\) with synthetic data
      \(\hat{\mathcal{D}}_g\). Training an algorithm \(A\) (e.g.\ a 5-layer
      perceptron) for task \(\mathcal{T}\) on
      \(\mathcal{D}^{\mathrm{train}}\) and \(\hat{\mathcal{D}}_g\) yields
      models \(M_r\) and \(\hat{M}_g\), respectively. Scoring both on the
      real test set \(\mathcal{D}^{\mathrm{test}}\) by a score \(S\)
      \footnote{WLOG, we assume a score on \(\mathbb{R}\) where higher is
      better} gives \(s_r\) and \(s_g\); the generator is useful for the
      pair \((\mathcal{T}, A)\) if \(s_g > s_r\), and the best generator is
      \(G_{g^{*}}\) with
      \(g^{*} \coloneq \max\{s_g : g \in \mathcal{G}\}\).

      \subsubsection{Augmented Training}
      \label{sec:augmented-training}
      We additionally train on the augmented set
      \(\tilde{\mathcal{D}}_j \coloneq \mathcal{D}^{\mathrm{train}} \cup
      \hat{\mathcal{D}}_j\), yielding \(\tilde{M}_j\), to assess whether
      synthetic data adds value on top of real training data.

    \subsubsection{Utility Evaluation Metrics}
    \label{sec:utility-evaluation}
      To make downstream evaluations comparable across tasks, we define a
      common toolbox of utility metrics used by researchers and
      practitioners. Given a downstream model \(M\) for task
      \(\mathcal{T}\), let \(\mathbf{\Pi}_M \coloneq (\pi_{t,i})_{t,i} \in
      \mathbb{R}^{L \times I}\) denote the trading portfolio generated by
      \(M\), where \(\pi_{t,i}\) is the trading order (position) in asset
      \(i\) at time \(t\). A utility metric \(U\) maps the portfolio to a
      scalar score \(u \coloneq U(\mathbf{\Pi}_M)\), instantiating the TSTR
      score \(S\) of §\ref{sec:utility-training-framework}, i.e.\
      \(s \equiv u = U(\mathbf{\Pi}_M)\) for
      \(M \in \{M_r, \hat{M}_j, \tilde{M}_j\}\). The per-period profit and
      loss (PnL) of the portfolio is
      \begin{equation}
        \mathrm{PnL}_t(\mathbf{\Pi}_M) \coloneq \sum_{i=1}^{I}
        \pi_{t,i}\, r_{t+1,i},
        \label{eq:utility-pnl}
      \end{equation}
      where \(r_{t+1,i}\) is the simple return of asset \(i\) over
      \([t, t+1)\). The toolbox evaluates the following metrics over the
      per-period PnL sequence
      \(\{\mathrm{PnL}_t(\mathbf{\Pi}_M)\}_{t=0}^{L-1}\):
      \begin{enumerate}[label=\textbf{[U\arabic*]}]
        \item \textbf{PnL.} Total profit and loss over the window,
        \(\mathrm{PnL}(\mathbf{\Pi}_M) \coloneq \sum_{t=0}^{L-1}
        \mathrm{PnL}_t(\mathbf{\Pi}_M)\).

        \item \textbf{Annualized Sharpe ratio.} The mean per-period PnL
        divided by its standard deviation,
        \begin{equation}
          \begin{split}
            \mathrm{Sharpe}(\mathbf{\Pi}_M) &\coloneq
            \frac{\hat{\mu}}{\hat{\sigma}}, \\\
            \hat{\mu} \coloneq \frac{1}{L}\sum_{t=0}^{L-1}
            \mathrm{PnL}_t, \quad &
            \hat{\sigma}^2 \coloneq \frac{1}{L-1}
            \sum_{t=0}^{L-1} \left(\mathrm{PnL}_t -
            \hat{\mu}\right)^2,
          \end{split}
          \label{eq:utility-sharpe}
        \end{equation}
        Since the evaluation horizon is a single trading year
        , no annualization factor (i.e.\ \(\sqrt{L}\))
        is applied.

        % \item \textbf{Entropic risk measure (ERM).} The risk-sensitive
        % certainty equivalent of the per-period PnL under exponential
        % utility with risk-aversion \(\gamma > 0\) \cite{Fecamp20},
        % \begin{equation}
        %   \mathrm{ERM}_\gamma(\mathbf{\Pi}_M) \coloneq
        %   -\frac{1}{\gamma}\log\frac{1}{L}\sum_{t=0}^{L-1}
        %   e^{-\gamma\, \mathrm{PnL}_t},
        %   \label{eq:utility-erm}
        % \end{equation}
        % which reduces to the mean PnL as \(\gamma \to 0\) and increasingly
        % penalizes losses as \(\gamma\) grows.

        \item \textbf{CVaR.} The conditional value at risk at level
        \(\alpha\), i.e.\ the expected loss in the worst \(\alpha\)-tail of
        the per-period PnL distribution \cite{Fecamp20},
        \begin{equation}
          \mathrm{CVaR}_\alpha(\mathbf{\Pi}_M) \coloneq
          -\mathbb{E}\left[\mathrm{PnL}_t \mid \mathrm{PnL}_t \leq
          \mathrm{VaR}_\alpha\right],
          \label{eq:utility-cvar}
        \end{equation}
        where \(\mathrm{VaR}_\alpha\) is the lower \(\alpha\)-quantile of
        the per-period PnL distribution.

        \item \textbf{Win rate.} The fraction of periods with positive
        PnL,
        \begin{equation}
          \mathrm{WR}(\mathbf{\Pi}_M) \coloneq \frac{1}{L}
          \sum_{t=0}^{L-1} \mathbb{I}\left\{\mathrm{PnL}_t >
          0\right\},
          \label{eq:utility-winrate}
        \end{equation}
        where \(\mathbb{I}\{\cdot\}\) is the indicator function.

        \item \textbf{Max drawdown.} The largest peak-to-trough decline
        of the cumulative PnL curve \(C_t \coloneq \sum_{s=0}^{t}
        \mathrm{PnL}_s\),
        \begin{equation}
          \mathrm{MDD}(\mathbf{\Pi}_M) \coloneq \max_{0 \leq u \leq v
          \leq L} \left(C_u - C_v\right).
          \label{eq:utility-mdd}
        \end{equation}
      \end{enumerate}

      Let \(u_n \coloneq U\left(\mathbf{\Pi}_M^{(n)}\right)\) be the value
      of a metric \(U\) on the \(n\)-th test window. We report each metric
      as
      \begin{equation}
          \bar{u} \coloneq \frac{1}{N}\sum_{n=1}^{N} u_n, \quad
          \hat{\sigma}_u \coloneq \sqrt{\frac{1}{N-1}
          \sum_{n=1}^{N} \left(u_n - \bar{u}\right)^2},
        \label{eq:utility-aggregation}
      \end{equation}
      the mean and standard deviation of \(U\) across all \(N \coloneq
      \vert\mathcal{D}^{\mathrm{test}}\vert\) test windows.

      Every task below instantiates the same protocol: train
      \(\{M_r, \hat{M}_j, \tilde{M}_j\}\) on
      \(\{\mathcal{D}^{\mathrm{train}}, \hat{\mathcal{D}}_j,
      \tilde{\mathcal{D}}_j\}\) per
      §\ref{sec:utility-training-framework}--\ref{sec:augmented-training},
      evaluate on \(\mathcal{D}^{\mathrm{test}}\), and report the toolbox
      metrics above.

\subsection{Option Delta Hedging}\label{sec:option-hedging}
    Option hedging is a popular downstream utility benchmark in market
    generator research \cite{Tanaka25, Boursin22}. We take the perspective
    of an options market maker selling a European call option on a cash
    equity to minimize counterparty risk when the call expires
    \cite{Fecamp20}.

    \subsubsection{Problem Definition}
    Consider a European call at strike \(K\), maturing at time \(L\), on
    cash equity \(i\) with payoff \(g(L) \coloneq (S_{L,i} - K)_+\),
    rebalanced at the discrete times \(\mathbb{T} \coloneq \{t \in
    \mathbb{N} : 0 \leq t \leq L-1\}\) (end-of-day). The hedger's policy
    is the hedge position \(\mathbf{\Delta} \coloneq (\Delta_t)_{t=0}^{L-1}
    \in \mathbb{R}^{L}\), i.e.\ the trading portfolio \(\mathbf{\Pi}_M
    \equiv \mathbf{\Delta}\). The loss is the quadratic replication
    error of the short call,
    \begin{equation}
      \mathcal{L}_{\Delta}(\mathbf{\Delta}) \coloneq
      \mathbb{E}\left[\left(g'(L) - \sum_{t=0}^{L-1} \Delta_{t}\,
      \tilde{g}(t)\, \left(e^{R_{t+1,i}} - 1\right)\right)^2\right],
      \label{eq:option-hedger-loss}
    \end{equation}
    where \(\tilde{g}(t) \coloneq S_{t,i}/K\) is the moneyness of the
    contract, \(\tilde{g}(L) \coloneq (S_{L,i}/K - 1)_+\) its payoff in
    moneyness units, and \(g'(L) \coloneq \tilde{g}(L) - c_0/K\) with
    \(c_0\) the call premium from the Black-Scholes-Merton equation. The
    optimal policy minimizes this loss,
    \begin{equation}
      \mathbf{\Delta}^{*} \coloneq \underset{\mathbf{\Delta}}{\arg\min}\;
      \mathcal{L}_{\Delta}(\mathbf{\Delta}).
      \label{eq:option-hedger-optimizer}
    \end{equation}

    \subsubsection{Framework}
    We use a long-short-term-memory (LSTM) network under the loss of
    equation \ref{eq:option-hedger-loss} \cite{Fecamp20}. While related
    work feeds raw price data \cite{Fecamp20,Boursin22,Tanaka25}, our
    framework recurrently feeds the moneyness \(\tilde{g}(t)\) instead of
    the raw price, allowing the network to generalize the payoff across
    strikes. At each step the LSTM takes the triplet
    \((\Delta_{t-1}, \tilde{g}(t), \tau_t)\) --- the previous hedge
    position, the current moneyness, and the normalized time-to-maturity
    \(\tau_t \coloneq 1 - t/L\).

    \subsubsection{Training}
    Since the LSTM hedges one asset at a time, each training window
    \(\mathbf{w} \in \mathcal{D}\) is treated as one channel. For each
    window we append an initial hedge \(\Delta(0) = 0\) and an initial
    moneyness \(\tilde{g}(0)\) \footnote{Equivalent to fixing the strike
    \(K = S_{0,i}/\tilde{g}(0)\) for the contract of maturity \(L\)},
    drawn from \(0.7\) to \(1.3\) at \(0.1\) increments \footnote{These
    values correspond to option contracts \(10\%\)--\(30\%\) out-of- and
    in-the-money relative to the initial price}, and propagate the
    log-moneyness recursively as
    \(\log \tilde{g}(t) = \log \tilde{g}(t-1) + R_{t,i}\)
    \footnote{For a training set of size \(N = \vert\mathcal{D}\vert\),
    this yields \(N \times 7\) training examples}. Models are trained per
    §\ref{sec:utility-training-framework}.

    \subsubsection{Evaluation}
    Models \(\{M_r, \hat{M}_j, \tilde{M}_j\}\) are scored on
    \(\mathcal{D}^{\mathrm{test}}\) with the toolbox of
    §\ref{sec:utility-evaluation}; we additionally report the hedging loss
    \(\mathcal{L}_{\Delta}\) of equation \ref{eq:option-hedger-loss},
    aggregated per equation \ref{eq:utility-aggregation}.

\subsection{Portfolio Hedging}
    Portfolio hedging mitigates the risk of adverse price movements in a
    market maker's inventory. Classical market-making models assume dealers
    hedge inventory risk externally once it becomes sufficiently large
    \cite{Gueant2011}; we instantiate this external hedge with \(J = 5\)
    liquid U.S.\ ETFs (table \ref{tab:asset-classification}) serving as
    proxies for market and sector exposures.

    \subsubsection{Problem Definition}
    For a time horizon \(L\) and a universe of \(I = 20\) inventory
    stocks, let \(P_t \coloneq (p_{t,i})_{i=1}^{I} \in \mathbb{R}^{I}\) be
    the market maker's inventory at time \(t\), modeled as continuous
    exposures. The hedge policy \(\mathbf{H} \coloneq (h_{t,j})_{t,j} \in
    \mathbb{R}^{L \times J}\) --- the trading portfolio \(\mathbf{\Pi}_M
    \coloneq \mathbf{H}\) --- nets the inventory notional exposure against
    the ETF hedge. The loss is the cumulative absolute net exposure,
    \begin{equation}
      \mathcal{L}_{H}(\mathbf{H}) \coloneq \sum_{t=1}^{L-1} \left\vert
      \sum_{i=1}^{I} R_{t+1,i}p_{t,i} - \sum_{j=1}^{J}
      R_{t+1,j}h_{t,j}\right\vert,
      \label{eq:portfolio-hedging-loss}
    \end{equation}
    where \(R_{t,i}\) and \(R_{t,j}\) are the daily log returns of the
    inventory stocks and ETFs, respectively. The optimal hedge minimizes
    this loss,
    \begin{equation}
      \mathbf{H}^{*} \coloneq \underset{\mathbf{H}}{\arg\min}\;
      \mathcal{L}_{H}(\mathbf{H}).
      \label{eq:portfolio-hedging-objective}
    \end{equation}

    \subsubsection{Inventory Simulation}
    Since proprietary inventory records are unavailable, a natural
    alternative is to simulate inventory with dedicated synthetic inventory
    generators \cite{Gueant2011}; however, this introduces external
    uncertainty that obscures the generator's utility in a controlled
    setting. We therefore evaluate on a static portfolio:
    \(P_t = \mathbf{1}, t \in [L]\), the \(I\)-dimensional unit vector.

    \subsubsection{Framework and Training}
    We apply the same LSTM framework as §\ref{sec:option-hedging}, feeding
    at each step the previous hedge \(\mathbf{H}_{t-1}\) and the current
    returns \(R_{t,i}, R_{t,j}\) for all \(i \in [I], j \in [J]\).
    Training follows §\ref{sec:utility-training-framework}.

    \subsubsection{Evaluation}
    Models \(\{M_r, \hat{M}_j, \tilde{M}_j\}\) are scored on
    \(\mathcal{D}^{\mathrm{test}}\) with the toolbox of
    §\ref{sec:utility-evaluation}, reported as the mean and standard
    deviation across all test windows.

\subsection{Alpha Generation}\label{sec:alpha-generation}
    Alpha generation takes the perspective of a hedge fund manager or
    proprietary trader seeking to capture excess returns by trading the
    full asset universe. Unlike the hedging tasks, which neutralize an
    existing exposure, the model takes directional positions to maximize
    risk-adjusted PnL.

    \subsubsection{Problem Definition}
    For a horizon \(L\) and the full universe of \(I = 25\) assets (20
    stocks and 5 ETFs, table \ref{tab:asset-classification}), the model
    \(M\) produces the trading portfolio \(\mathbf{\Pi}_M \coloneq
    (\pi_{t,i})_{t,i} \in \mathbb{R}^{L \times I}\), where \(\pi_{t,i}\) is
    the position in asset \(i\) held over \([t, t+1)\), with per-period PnL
    as in equation \ref{eq:utility-pnl}. The loss is the negative
    differentiable Sharpe ratio,
    \begin{equation}
      \mathcal{L}_{\Pi}(\mathbf{\Pi}_M) \coloneq -
      \mathrm{Sharpe}(\mathbf{\Pi}_M) = -\frac{\hat{\mu}}{\hat{\sigma}},
      \label{eq:alpha-sharpe-loss}
    \end{equation}
    where \(\mathrm{Sharpe}(\cdot)\) is defined in equation
    \ref{eq:utility-sharpe} (and equals the annualized value at the
    one-trading-year horizon). The optimal portfolio is
    \begin{equation}
        \mathbf{\Pi}^{*} \coloneq
        \underset{\mathbf{\Pi}}{\arg\min}\;
        \mathcal{L}_{\Pi}(\mathbf{\Pi}) 
        = \underset{\mathbf{\Pi}}{\arg\max}\;
        \mathrm{Sharpe}(\mathbf{\Pi}).
      \label{eq:alpha-objective}
    \end{equation}

    \subsubsection{Framework}
    We employ the same LSTM architecture as §\ref{sec:option-hedging}, with
    an output head emitting cross-sectional trading orders for all \(I\)
    assets. At each step \(t\), the LSTM feeds the previous order vector
    \(\mathbf{\Pi}_{t-1}\) and the current log returns \(\mathbf{R}_t\), and
    outputs the updated orders \(\mathbf{\Pi}_t \coloneq
    (\pi_{t,i})_{i=1}^{I} \in \mathbb{R}^{I}\). Training follows
    §\ref{sec:utility-training-framework}.

    \subsubsection{Evaluation}
    Models \(\{M_r, \hat{M}_j, \tilde{M}_j\}\) are scored on
    \(\mathcal{D}^{\mathrm{test}}\) with the toolbox of
    §\ref{sec:utility-evaluation}  reported as the mean and standard deviation across all
    \(N\) test windows.
